\begin{proof}[Proof of \cref{obj:lem:realized-dose-mean}]
\begin{enumerate}
\item Define the projection onto the dose interval and the normalized stratum laws by
\[
\pi_{[0,1]}:\mathbb R\longrightarrow[0,1],
\qquad
\pi_{[0,1]}(a)=\max\{0,\min\{1,a\}\},
\]
\[
p_x=P(X=x),
\qquad
P_x=\frac{1}{p_x}P|_{\{X=x\}},
\qquad
\gamma_x=\mathcal L_{P_x}(A),
\qquad
\lambda_x=\mathcal L_{P_x}(Y(\cdot)),
\qquad x\in\mathcal X.
\]
By \cref{obj:ass:stratum-overlap}, \(p_x\ge p_{\min}>0\), so \(P_x\), \(\gamma_x\), and \(\lambda_x\) are probability measures. The law \(P_x\) is concentrated on \(\{X=x\}\).

Fix \(x\in\mathcal X\). For measurable sets
\[
\mathcal C_{\mathcal S}\in\mathcal E,
\qquad
\mathcal C_A\in\mathcal B(\mathbb R),
\]
\cref{obj:ass:conditional-unconfoundedness} gives
\[
\begin{aligned}
P_x\{Y(\cdot)\in\mathcal C_{\mathcal S},\,A\in\mathcal C_A\}
&=\frac{P\{Y(\cdot)\in\mathcal C_{\mathcal S},\,A\in\mathcal C_A,\,X=x\}}{p_x}\\
&=\frac{P\{Y(\cdot)\in\mathcal C_{\mathcal S},\,X=x\}}{p_x}
  \frac{P\{A\in\mathcal C_A,\,X=x\}}{p_x}\\
&=\lambda_x(\mathcal C_{\mathcal S})\gamma_x(\mathcal C_A).
\end{aligned}
\]
Thus equality on measurable rectangles yields
\[
\mathcal L_{P_x}(A,Y(\cdot))=\gamma_x\otimes\lambda_x.
\]
For every \(a\in\mathbb R\), \cref{obj:ass:bounded-potential-outcomes} transfers to the normalized restriction and gives
\[
\lambda_x\bigl\{f\in\mathcal S:
e(f,\pi_{[0,1]}(a))\notin[0,1]\bigr\}=0.
\]
The evaluation map from \cref{obj:def:potential-path-space} is jointly measurable. Integrating these fixed-dose assertions under the product law therefore gives
\[
\int_{\mathbb R}\int_{\mathcal S}
\mathbf1\{e(f,\pi_{[0,1]}(a))\notin[0,1]\}
\,\lambda_x(df)\,\gamma_x(da)=0,
\]
and hence
\[
P_x\{e(Y(\cdot),\pi_{[0,1]}(A))\in[0,1]\}=1.
\]
By \cref{obj:ass:latent-support}, \(\pi_{[0,1]}(A)=A\) almost surely under \(P_x\). Consistency from \cref{obj:ass:consistency} consequently yields
\[
P_x\{Y\in[0,1]\}=1.
\]
Combining the two strata,
\[
P\{Y\notin[0,1]\}
=\sum_{x\in\mathcal X}p_xP_x\{Y\notin[0,1]\}=0.
\]
% lean: CausalSmith.Stat.NoisydoseWeakdesignTransition.realized_outcome_mem_Icc

\item Define a mean extension on the entire real dose axis by
\[
\widetilde\mu_x(a)=\mu_x(\pi_{[0,1]}(a)),
\qquad (a,x)\in\mathbb R\times\mathcal X.
\]
The parameter restriction gives \(\beta>0\), and
\cref{obj:ass:holder-mean} supplies
\[
|\mu_x(a)-\mu_x(a')|\le |a-a'|^\beta,
\qquad a,a'\in[0,1],\quad x\in\mathcal X.
\]
Thus each \(\mu_x\) is continuous on \([0,1]\). Composition with the continuous projection, together with the discrete stratum space, makes
\((a,x)\mapsto\widetilde\mu_x(a)\) measurable.
By \cref{obj:ass:mean-range},
\[
0\le\widetilde\mu_x(a)\le1,
\qquad (a,x)\in\mathbb R\times\mathcal X.
\]
In particular,
\[
|Y|\le1\quad P\text{-almost surely},
\qquad
|\widetilde\mu_X(A)|\le1.
\]
Both random variables are measurable and integrable under the probability law \(P\).
Finally, latent support gives
\[
\widetilde\mu_X(A)=\mu_X(A)
\qquad P\text{-almost surely}.
\]
% lean: CausalSmith.Stat.NoisydoseWeakdesignTransition.projectedMean_measurable
% lean: CausalSmith.Stat.NoisydoseWeakdesignTransition.realized_outcome_integrable
% lean: CausalSmith.Stat.NoisydoseWeakdesignTransition.projectedMean_random_integrable
% lean: CausalSmith.Stat.NoisydoseWeakdesignTransition.realized_dose_mean (hactual)

\item Fix a bounded measurable test function
\[
F:\mathbb R\times\mathcal X\times\mathbb R\longrightarrow\mathbb R,
\qquad
|F(a,x,z)|\le B_F,
\qquad B_F\in[0,\infty).
\]
The preceding bounds imply
\[
|F(A,X,Z)Y|
=|F(A,X,Z)|\,|Y|\le B_F
\quad P\text{-almost surely},
\]
\[
|F(A,X,Z)\widetilde\mu_X(A)|
=|F(A,X,Z)|\,|\widetilde\mu_X(A)|\le B_F.
\]
These measurable products are integrable. Splitting the first expectation over the two strata gives
\[
\begin{aligned}
E_P[F(A,X,Z)Y]
&=\sum_{x\in\mathcal X}
 E_P[\mathbf1\{X=x\}F(A,X,Z)Y]\\
&=\sum_{x\in\mathcal X}p_x
 \int F(A,x,Z)Y\,dP_x.
\end{aligned}
\]
Here the second equality follows from the definition of \(P_x\) and \(X=x\) almost surely under \(P_x\). The same calculation gives
\[
E_P[F(A,X,Z)\widetilde\mu_X(A)]
=\sum_{x\in\mathcal X}p_x
 \int F(A,x,Z)\widetilde\mu_x(A)\,dP_x.
\]
% lean: CausalSmith.Stat.NoisydoseWeakdesignTransition.integral_eq_stratum_sum
% lean: CausalSmith.Stat.NoisydoseWeakdesignTransition.bounded_test_mean

\item Fix a stratum \(x\). Define the error law, clipped evaluation, and product integrand by
\[
\zeta=\mathcal L_P(Z),
\]
\[
e_{\mathrm{bd}}:\mathcal S\times\mathbb R\longrightarrow[0,1],
\qquad
e_{\mathrm{bd}}(f,a)
=\max\{0,\min\{1,e(f,\pi_{[0,1]}(a))\}\},
\]
\[
G_{x,F}:\mathbb R\times\mathbb R\times\mathcal S\longrightarrow\mathbb R,
\qquad
G_{x,F}(z,a,f)=F(a,x,z)e_{\mathrm{bd}}(f,a).
\]
These functions are measurable, and
\[
0\le e_{\mathrm{bd}}(f,a)\le1,
\qquad
|G_{x,F}(z,a,f)|
\le |F(a,x,z)|\le B_F.
\]
Thus \(G_{x,F}\) is integrable under \(\zeta\otimes(\gamma_x\otimes\lambda_x)\).

Error independence from \cref{obj:ass:error-independence} survives restriction to the positive-probability stratum. Indeed, for
\[
\mathcal C_Z\in\mathcal B(\mathbb R),
\qquad
\mathcal C_{A,\mathcal S}\in\mathcal B(\mathbb R)\otimes\mathcal E,
\]
\[
\begin{aligned}
&P_x\{Z\in\mathcal C_Z,\,(A,Y(\cdot))\in\mathcal C_{A,\mathcal S}\}\\
&\quad=
\frac{P\{Z\in\mathcal C_Z,\,(A,Y(\cdot))\in\mathcal C_{A,\mathcal S},\,X=x\}}{p_x}\\
&\quad=
\zeta(\mathcal C_Z)
\frac{P\{(A,Y(\cdot))\in\mathcal C_{A,\mathcal S},\,X=x\}}{p_x}\\
&\quad=
\zeta(\mathcal C_Z)P_x\{(A,Y(\cdot))\in\mathcal C_{A,\mathcal S}\}.
\end{aligned}
\]
Equality on rectangles and the schedule--dose product law already established yield
\[
\mathcal L_{P_x}(Z,A,Y(\cdot))
=\zeta\otimes(\gamma_x\otimes\lambda_x).
\]
The random-evaluation bound, latent support, and consistency also give
\[
e_{\mathrm{bd}}(Y(\cdot),A)=Y
\qquad P_x\text{-almost surely}.
\]
Consequently, change of variables under this product law and Fubini's theorem give
\[
\int F(A,x,Z)Y\,dP_x
=
\int_{\mathbb R}\int_{\mathbb R}\int_{\mathcal S}
G_{x,F}(z,a,f)\,\lambda_x(df)\,\gamma_x(da)\,\zeta(dz).
\]

For every \(a\in\mathbb R\), the fixed-dose bound established above makes clipping leave the evaluation unchanged \(\lambda_x\)-almost everywhere. The normalized restriction formula and
\cref{obj:def:conditional-potential-mean} therefore give
\[
\begin{aligned}
\int_{\mathcal S}e_{\mathrm{bd}}(f,a)\,\lambda_x(df)
&=\int_{\mathcal S}e(f,\pi_{[0,1]}(a))\,\lambda_x(df)\\
&=\frac{E_P[\mathbf1\{X=x\}Y(\pi_{[0,1]}(a))]}{p_x}\\
&=\mu_x(\pi_{[0,1]}(a))
=\widetilde\mu_x(a).
\end{aligned}
\]
Integrating out the schedule thus gives
\[
\int F(A,x,Z)Y\,dP_x
=
\int_{\mathbb R}\int_{\mathbb R}
F(a,x,z)\widetilde\mu_x(a)\,\gamma_x(da)\,\zeta(dz).
\]
Moreover,
\[
|F(a,x,z)\widetilde\mu_x(a)|\le B_F,
\qquad a,z\in\mathbb R,
\]
so this measurable integrand is integrable. Marginalizing the product law, using
\(\lambda_x(\mathcal S)=1\), gives
\[
\mathcal L_{P_x}(Z,A)=\zeta\otimes\gamma_x.
\]
It follows that
\[
\int F(A,x,Z)Y\,dP_x
=
\int F(A,x,Z)\widetilde\mu_x(A)\,dP_x.
\]
% lean: CausalSmith.Stat.NoisydoseWeakdesignTransition.bounded_test_mean_stratum

Substituting this equality into the two stratum sums proves, for every bounded measurable \(F\),
\[
E_P[F(A,X,Z)Y]
=
E_P[F(A,X,Z)\widetilde\mu_X(A)].
\]
% lean: CausalSmith.Stat.NoisydoseWeakdesignTransition.bounded_test_mean

\item Apply this identity to indicator tests. For every
\[
\mathcal C_{AXZ}\in\mathcal B(\mathbb R\times\mathcal X\times\mathbb R),
\]
it gives
\[
\int_{\{(A,X,Z)\in\mathcal C_{AXZ}\}}Y\,dP
=
\int_{\{(A,X,Z)\in\mathcal C_{AXZ}\}}\widetilde\mu_X(A)\,dP.
\]
Every event in \(\sigma(A,X,Z)\) has this preimage form. The candidate
\(\widetilde\mu_X(A)\) is integrable and \(\sigma(A,X,Z)\)-measurable, so the defining set-integral characterization and uniqueness of conditional expectation yield
\[
E_P[Y\mid A,X,Z]=\widetilde\mu_X(A)
=\mu_X(A)
\qquad P\text{-almost surely},
\]
where the second equality uses latent support.
% lean: CausalSmith.Stat.NoisydoseWeakdesignTransition.condExp_eq_of_bounded_tests
% lean: CausalSmith.Stat.NoisydoseWeakdesignTransition.realized_dose_mean (hXZ)

\item For any bounded measurable function
\[
F_{AX}:\mathbb R\times\mathcal X\longrightarrow\mathbb R,
\]
define its extension by
\[
F_{AXZ}:\mathbb R\times\mathcal X\times\mathbb R\longrightarrow\mathbb R,
\qquad
F_{AXZ}(a,x,z)=F_{AX}(a,x).
\]
This extension is measurable and has the same bound. The bounded-test identity therefore supplies
\[
E_P[F_{AX}(A,X)Y]
=
E_P[F_{AX}(A,X)\widetilde\mu_X(A)].
\]
The same indicator-test characterization, now for \(\sigma(A,X)\), gives
\[
E_P[Y\mid A,X]=\widetilde\mu_X(A)
=\mu_X(A)
\qquad P\text{-almost surely}.
\]
Finally, replacing the projected mean by its almost-surely equal value in the original bounded-test identity gives
\[
E_P[F(A,X,Z)Y]
=
E_P[F(A,X,Z)\widetilde\mu_X(A)]
=
E_P[F(A,X,Z)\mu_X(A)].
\]
% lean: CausalSmith.Stat.NoisydoseWeakdesignTransition.condExp_eq_of_bounded_tests
% lean: CausalSmith.Stat.NoisydoseWeakdesignTransition.realized_dose_mean
\end{enumerate}
\end{proof}
